\documentclass[12pt,a4paper]{article}

\usepackage{fontspec}
\setmainfont{Times New Roman}
\usepackage[english,russian]{babel}
\usepackage{amsmath,amssymb}
\usepackage{bm}
\usepackage{booktabs}
\usepackage{graphicx}
\usepackage{geometry}
\usepackage[hidelinks]{hyperref}

\geometry{left=30mm,right=15mm,top=20mm,bottom=20mm}
\graphicspath{{figures/}}
\numberwithin{equation}{section}

\newcommand{\studentName}{Мозжевилова Екатерина}
\newcommand{\studentGroup}{5030102/30201}

\newsavebox{\criterioncontent}
\newcommand{\theorypart}[1]{%
    \par\medskip
    \noindent\textbf{#1}\par\nopagebreak\smallskip
}

\newenvironment{criterionbox}
{%
    \par\medskip
    \begin{center}
    \begin{lrbox}{\criterioncontent}
    \begin{minipage}{0.90\textwidth}
    \textbf{Теорема (критерий потери полного столбцового ранга).}\par\smallskip
}
{%
    \end{minipage}
    \end{lrbox}
    \setlength{\fboxsep}{9pt}%
    \fbox{\usebox{\criterioncontent}}
    \end{center}
    \medskip
}

\begin{document}

\begin{titlepage}
    \centering
    {\large Санкт-Петербургский политехнический университет Петра Великого\par}
    \vfill
    {\Large\bfseries Интервальный Анализ\par}
    \vspace{0.5cm}
    {\large Лабораторная № 1\par}
    \vfill
    \begin{flushright}
        Выполнила: \studentName\\
        Группа: \studentGroup
    \end{flushright}
    \vfill
    Санкт-Петербург\\
    2026
\end{titlepage}

\section{Постановка задачи}

\begin{enumerate}
    \setlength{\itemsep}{0.8\baselineskip}

    \item Для матрицы $\boldsymbol{A}_1\in\mathbb{IR}^{3\times2}$
    \[
        \operatorname{mid}\boldsymbol{A}_1=
        \begin{pmatrix}
            0.95 & 1.00 \\
            1.05 & 1.00 \\
            1.10 & 1.00
        \end{pmatrix}
    \]
    заданы радиусы:
    \[
        \operatorname{rad}\boldsymbol{A}_1=\delta
        \begin{pmatrix}
            1 & 1 \\
            1 & 1 \\
            1 & 1
        \end{pmatrix},
        \qquad\text{или}\qquad
        \operatorname{rad}\boldsymbol{A}_1=\delta
        \begin{pmatrix}
            1 & 0 \\
            1 & 0 \\
            1 & 0
        \end{pmatrix},
        \qquad \delta\geq0.
    \]
    Первый вариант рассматривается для задачи малоракурсной томографии,
    второй --- для линейной регрессии.

    \item Матрица $\boldsymbol{A}_2\in\mathbb{IR}^{3\times3}$
    \[
        \operatorname{mid}\boldsymbol{A}_2=
        \begin{pmatrix}
            1.10 & 0.90 & 1.10 \\
            1.40 & 1.00 & 0.80 \\
            0.80 & 1.40 & 1.20
        \end{pmatrix},
        \qquad
        \operatorname{rad}\boldsymbol{A}_2=\delta
        \begin{pmatrix}
            1 & 1 & 1 \\
            1 & 1 & 1 \\
            1 & 1 & 1
        \end{pmatrix},
        \qquad \delta\geq0.
    \]
\end{enumerate}

Для всех заданных вариантов матриц радиусов требуется:
\begin{enumerate}
    \item найти диапазоны $\delta$, при которых $\boldsymbol{A}_1$ и
    $\boldsymbol{A}_2$ содержат матрицы неполного столбцового ранга;
    \item при минимальном значении $\delta$ построить точечную матрицу
    $A'_k\in\boldsymbol{A}_k$ и ненулевой вектор $x_*$ такие, что
    $A'_kx_*=0$, где $k\in\{1,2\}$;
    \item проверить принадлежность и ранг построенных матриц и обсудить
    полученные результаты для линейной регрессии и малоракурсной томографии.
\end{enumerate}

\section{Теоретическая часть}

\subsection{Критерий потери полного столбцового ранга}

Пусть $A_c=(a_{ij})$ и $D=(d_{ij})$ --- точечные матрицы из
$\mathbb{R}^{m\times n}$, $m\geq n$, причём $D\geq0$. Здесь
$i=1,\ldots,m$, $j=1,\ldots,n$, а неравенства между матрицами и
векторами понимаются покомпонентно.

Рассмотрим интервальную матрицу
\[
    \boldsymbol{A}(\delta)
    = [A_c-\delta D,\,A_c+\delta D],
    \qquad \delta\geq0,
\]
где $A_c=\operatorname{mid}\boldsymbol{A}$ и
$\delta D=\operatorname{rad}\boldsymbol{A}$.

\begin{criterionbox}
Матрица $\boldsymbol{A}(\delta)$ содержит точечную матрицу неполного
столбцового ранга тогда и только тогда, когда существует вектор
$x\neq0$, для которого
\[
    |A_cx|\leq\delta D|x|.
\]
Модуль вектора и матрицы вычисляется покомпонентно.
\end{criterionbox}

Этот критерий приведён в пособии~[1]. Для прямоугольной матрицы он
характеризует именно потерю полного столбцового ранга.

\theorypart{Следствие}
Минимальный радиус равен
\[
    \delta_* = \min_{x\neq0}\max_i
    \frac{|(A_cx)_i|}{(D|x|)_i}.
\]
Отношение с нулевым знаменателем считается допустимым только при
нулевом числителе. Поскольку семейства $\boldsymbol{A}(\delta)$
вложены, искомый диапазон имеет вид
\[
    \delta\in[\delta_*,+\infty).
\]

\subsection{Задача линейной регрессии}

\theorypart{Исходная матрица}
Для регрессии с одним признаком и свободным коэффициентом
\[
    A_c=
    \begin{pmatrix}
        t_1 & 1\\
        \vdots & \vdots\\
        t_m & 1
    \end{pmatrix},
    \qquad
    D=
    \begin{pmatrix}
        1 & 0\\
        \vdots & \vdots\\
        1 & 0
    \end{pmatrix}.
\]
Значения $t_i$ считаются неточными, а столбец единиц --- точным.

\theorypart{Упрощение критерия}
Для $x=(x_1,x_2)^T$ критерий принимает вид
\[
    |t_ix_1+x_2|\leq\delta|x_1|,
    \qquad i=1,\ldots,m.
\]
Здесь $x_1\neq0$. Введя $c=-x_2/x_1$, получаем
\[
    |t_i-c|\leq\delta.
\]
Следовательно, необходимо выбрать $c$ так, чтобы наибольшее из
расстояний $|t_i-c|$ было минимальным:
\[
    \delta_* = \min_{c\in\mathbb{R}}\max_i |t_i-c|.
\]

\theorypart{Критический радиус}
Наибольшее расстояние определяется крайними значениями $t_i$.
В оптимальной точке $c_*$ расстояния до концов диапазона равны:
\[
    c_* - \min_i t_i = \max_i t_i - c_*.
\]
Отсюда
\[
    c_* = \frac{\min_i t_i+\max_i t_i}{2},
    \qquad
    \delta_* = \frac{\max_i t_i-\min_i t_i}{2}.
\]

\begin{samepage}
\theorypart{Матрица на границе}
При $\delta=\delta_*$ для каждого $i$ выполняется
\[
    |t_i-c_*|\leq\delta_*,
    \qquad
    c_*\in[t_i-\delta_*,\,t_i+\delta_*].
\]
Следовательно, в каждой строке первый элемент можно выбрать равным
$c_*$, а второй элемент остаётся равным единице. Получаем матрицу
\[
    A'=\begin{pmatrix}
        c_* & 1\\
        \vdots & \vdots\\
        c_* & 1
    \end{pmatrix},
    \qquad
    x_*=\begin{pmatrix}1\\-c_*\end{pmatrix},
    \qquad A'x_*=0.
\]
\end{samepage}

\subsection{Задача малоракурсной томографии}

\theorypart{Упрощение критерия}
В рассматриваемом варианте все элементы матрицы $D$ равны единице.
Поэтому для каждой строки
\[
    (D|x|)_i=\sum_j|x_j|=\|x\|_1.
\]
Также
\[
    \max_i |(A_cx)_i|=\|A_cx\|_\infty.
\]
Тогда критерий потери ранга принимает вид
\[
    \delta_* = \min_{x\neq0}
    \frac{\|A_cx\|_\infty}{\|x\|_1}
    =\min_{\|x\|_1=1}\|A_cx\|_\infty.
\]

\theorypart{Переход к линейному программированию}
Зафиксируем знаки и запишем
$x_j=s_ju_j$, где $s_j\in\{-1,1\}$ и $u_j\geq0$. Для каждого набора
знаков получаем задачу линейного программирования
\[
    \begin{gathered}
        \min_{u,\delta}\ \delta,\\[0.4em]
        \left\{
        \begin{aligned}
            -\delta&\leq\sum_{j=1}^{n}a_{ij}s_ju_j\leq\delta,
                && i=1,\ldots,m,\\
            u_j&\geq0,
                && j=1,\ldots,n,\\
            \sum_{j=1}^{n}u_j&=1.
        \end{aligned}
        \right.
    \end{gathered}
\]

Наименьшее значение по всем наборам знаков даёт $\delta_*$, а
соответствующее решение определяет
$x_{*j}=s_ju_{*j}$.

\theorypart{Матрица на границе}
Искомую матрицу представим в виде $A'=A_c+E$. Чтобы выполнялись
$A'\in\boldsymbol{A}(\delta_*)$ и $A'x_*=0$, достаточно построить
$E$ так, чтобы
\[
    |E|\leq\delta_*D,
    \qquad
    Ex_*=-A_cx_*.
\]
Для найденного вектора $x_{*j}=s_ju_{*j}$ положим
\[
    e_{ij}=-(A_cx_*)_i s_j.
\]
Из ограничений линейной задачи следует
\[
    |(A_cx_*)_i|\leq\delta_*,
\]
поэтому $|E|\leq\delta_*D$. Так как $\sum_j u_{*j}=1$, имеем
\[
    (Ex_*)_i
    =-(A_cx_*)_i\sum_{j=1}^{n}s_j^2u_{*j}
    =-(A_cx_*)_i.
\]
Оба требуемых условия выполнены. Следовательно,
\[
    A'=A_c+E\in\boldsymbol{A}(\delta_*)
\]
удовлетворяет равенству $A'x_*=0$ и имеет неполный столбцовый ранг.

\section{Реализация}

\subsection{Средства реализации}

\noindent
Python --- основной язык программы.\\
NumPy --- операции с векторами и матрицами.\\
SciPy --- решение задач линейного программирования функцией
\texttt{linprog} с методом \texttt{highs}.

\subsection{Структура программы}

\theorypart{Основные функции}
\noindent
\texttt{solve\_regression(center)} --- находит $\delta_*$, $x_*$ и $A'$
для задачи линейной регрессии.\par
\smallskip
\noindent
\texttt{solve\_tomography(center)} --- находит $\delta_*$, $x_*$ и $A'$
для задачи малоракурсной томографии.\par
\smallskip
\noindent
\texttt{build\_tomography\_matrix(center, D, x)} --- строит матрицу $A'$
для задачи малоракурсной томографии по найденному вектору $x_*$.

\theorypart{Вспомогательные функции}
\noindent
\texttt{maximal\_minors(matrix)} --- вычисляет максимальные миноры
прямоугольной матрицы.\par
\smallskip
\noindent
\texttt{print\_result(...)} --- печатает найденные $\delta_*$, $x_*$ и $A'$,
а также результаты проверок.

\subsection{solve\_regression}

\noindent
\textbf{Входные данные:} \texttt{center} --- центральная матрица
$A_c$. Пусть $m$ --- число её строк. Тогда
$A_c\in\mathbb{R}^{m\times2}$, причём второй столбец состоит из единиц.

\noindent
\textbf{Выходные данные:} кортеж \texttt{(delta, x, A\_prime)}, содержащий
критический радиус $\delta_*$, вектор $x_*$ и матрицу $A'$.

\medskip
\noindent
Функция реализует упрощённый алгоритм, полученный в теоретической части
для регрессионной задачи.

\smallskip
\noindent
\textbf{Обозначения:}\\
\texttt{center} --- центральная матрица $A_c$;\\
\texttt{values} --- значения $t_i$ из первого столбца;\\
\texttt{minimum}, \texttt{maximum} --- минимальное и максимальное
значения $t_i$;\\
\texttt{new\_column\_value} --- величина $c_*$;\\
\texttt{delta} --- критический радиус $\delta_*$;\\
\texttt{x} --- вектор $x_*$;\\
\texttt{A\_prime} --- матрица $A'$.

\theorypart{Псевдокод}

\begin{enumerate}
    \setlength{\itemsep}{0.3\baselineskip}
    \item Если \texttt{center} не имеет двух столбцов или её второй
    столбец не состоит из единиц, завершить работу с ошибкой.
    \item \texttt{values} $\gets$ первый столбец \texttt{center}.
    \item \texttt{minimum} $\gets$ \texttt{min(values)},
    \texttt{maximum} $\gets$ \texttt{max(values)}.
    \item \texttt{delta} $\gets$
    \texttt{(maximum - minimum) / 2}.
    \item \texttt{new\_column\_value} $\gets$
    \texttt{(minimum + maximum) / 2}.
    \item \texttt{A\_prime} $\gets$ копия \texttt{center};
    первый столбец \texttt{A\_prime} $\gets$
    \texttt{new\_column\_value}.
    \item \texttt{x} $\gets$ \texttt{[1, -new\_column\_value]}.
    \item Вернуть \texttt{delta}, \texttt{x}, \texttt{A\_prime}.
\end{enumerate}

\theorypart{Сложность}

\noindent
Обозначим число строк входной матрицы $A_c$ через $m$.

\smallskip
\noindent
\textbf{Временная сложность:} $O(m)$. Проверка данных, поиск крайних
значений и копирование матрицы требуют одного прохода по её строкам.

\smallskip
\noindent
\textbf{Затраты памяти:} $O(m)$. Память требуется для возвращаемой
матрицы $A'$ размера $m\times2$.

\subsection{solve\_tomography}

\noindent
\textbf{Входные данные:} \texttt{center} --- центральная матрица $A_c$.
Пусть $m$ --- число её строк, $n$ --- число столбцов. Тогда
$A_c\in\mathbb{R}^{m\times n}$.

\noindent
\textbf{Выходные данные:} кортеж \texttt{(delta, x, A\_prime)}, содержащий
критический радиус $\delta_*$, вектор $x_*$ и матрицу $A'$.

\medskip
\noindent
Функция реализует минимаксный алгоритм из теоретической части для задачи
малоракурсной томографии.

\smallskip
\noindent
\textbf{Обозначения:}\\
\texttt{center} --- центральная матрица $A_c$;\\
\texttt{row\_count}, \texttt{column\_count} --- числа строк $m$
и столбцов $n$;\\
\texttt{signs} --- набор знаков $s$;\\
\texttt{u} --- вектор модулей компонент $x$;\\
\texttt{signed\_center} --- матрица коэффициентов после подстановки
$x_j=s_ju_j$;\\
\texttt{best\_delta} --- критический радиус $\delta_*$;\\
\texttt{best\_x} --- вектор $x_*$;\\
\texttt{A\_prime} --- матрица $A'$.

\theorypart{Псевдокод}

\begin{enumerate}
    \setlength{\itemsep}{0.3\baselineskip}
    \item Если \texttt{center} не является двумерной матрицей,
    завершить работу с ошибкой.
    \item \texttt{best\_delta} $\gets +\infty$,
    \texttt{best\_x} $\gets$ неопределённое значение.
    \item Для каждого \texttt{signs}$\in\{-1,1\}^n$:
    \begin{enumerate}
        \item \texttt{signed\_center} $\gets$ столбцы \texttt{center},
        умноженные на соответствующие знаки;
        \item решить с помощью \texttt{linprog} задачу
        \[
            \min\delta:\quad
            -\delta\leq A_c\operatorname{diag}(s)u\leq\delta,
            \quad u\geq0,\quad \sum_j u_j=1;
        \]
        \item если найденное значение меньше \texttt{best\_delta},\\
        сохранить его и восстановить
        \texttt{best\_x} $\gets s\circ u$.
    \end{enumerate}
    \item Если \texttt{best\_x} не найден, завершить работу с ошибкой.
    \item $D\gets$ матрица единиц размера $m\times n$.
    \item \texttt{A\_prime} $\gets$
    \texttt{build\_tomography\_matrix(center, D, best\_x)}.
    \item Вернуть \texttt{best\_delta}, \texttt{best\_x},
    \texttt{A\_prime}.
\end{enumerate}

\theorypart{Сложность}

\noindent
Здесь $m$ --- число строк, $n$ --- число столбцов входной матрицы $A_c$.

\smallskip
\noindent
\textbf{Число задач линейного программирования:} $2^n$.
Каждая задача содержит $n+1$ переменных, $2m$ неравенств и одно
равенство.

\smallskip
\noindent
\textbf{Время подготовки задач:} $O(2^nmn)$.

\noindent
Для каждого набора знаков формируются матрицы ограничений размера,
пропорционального $mn$.

\smallskip
\noindent
Обозначим через $T_{\mathrm{LP}}(m,n)$ время решения одной задачи
линейного программирования с $n+1$ переменными, $2m$ неравенствами
и одним равенством.

\smallskip
\noindent
\textbf{Временная сложность:}
$O\!\left(2^n\bigl(mn+T_{\mathrm{LP}}(m,n)\bigr)\right)$.

\smallskip
\noindent
\textbf{Затраты памяти:} $O(mn)$ без учёта внутренней памяти
\texttt{linprog}. Одновременно хранится только одна линейная задача.

\subsection{build\_tomography\_matrix}

\noindent
\textbf{Входные данные:} \texttt{center} --- матрица $A_c$,
\texttt{D} --- неотрицательная матрица радиусов того же размера,
\texttt{x} --- найденный вектор $x_*$. Пусть $m$ и $n$ --- числа строк
и столбцов этих матриц; тогда $A_c,D\in\mathbb{R}^{m\times n}$,
$x_*\in\mathbb{R}^n$.

\noindent
\textbf{Выходные данные:} матрица \texttt{A\_prime}, соответствующая
$A'$ и удовлетворяющая равенству $A'x_*=0$.

\medskip
\noindent
Функция реализует построение матрицы на границе интервального семейства,
описанное в теоретической части для задачи томографии.

\smallskip
\noindent
\textbf{Обозначения:}\\
\texttt{center} --- центральная матрица $A_c$;\\
\texttt{D} --- матрица радиусов $D$;\\
\texttt{x} --- вектор $x_*$;\\
\texttt{residual} --- остаток $r=A_cx_*$;\\
\texttt{capacity} --- допустимое изменение $p=D|x_*|$;\\
\texttt{row\_factors} --- коэффициенты $q_i$;\\
\texttt{E} --- возмущение $E$;\\
\texttt{A\_prime} --- матрица $A'$.

\theorypart{Псевдокод}

\begin{enumerate}
    \setlength{\itemsep}{0.3\baselineskip}
    \item \texttt{residual} $\gets A_cx_*$.
    \item \texttt{capacity} $\gets D|x_*|$.
    \item \texttt{row\_factors} $\gets$ нулевой вектор длины $m$.
    \item Для каждой строки $i$:
    \begin{enumerate}
        \item если $p_i=0$ и $r_i\neq0$, завершить работу с ошибкой;
        \item если $p_i>0$, положить $q_i\gets-r_i/p_i$.
    \end{enumerate}
    \item $e_{ij}\gets q_i d_{ij}\operatorname{sign}(x_{*j})$.
    \item \texttt{A\_prime} $\gets A_c+E$.
    \item Вернуть \texttt{A\_prime}.
\end{enumerate}

\theorypart{Сложность}

\noindent
Здесь $m$ и $n$ --- числа строк и столбцов входных матриц.

\smallskip
\noindent
\textbf{Временная сложность:} $O(mn)$. Умножения матриц на вектор
и построение всех элементов $E$ требуют обработки $mn$ элементов.

\smallskip
\noindent
\textbf{Затраты памяти:} $O(mn)$. Основную память занимают матрицы
$E$ и $A'$; три вспомогательных вектора занимают $O(m)$ памяти.

\section{Результаты}

Для каждого варианта программа нашла критический радиус $\delta_*$,
ненулевой вектор $x_*$ и матрицу $A'$ неполного столбцового ранга.
Результаты проверок приведены в таблице~\ref{tab:results-check}.

\begin{table}[ht]
    \centering
    \caption{Проверка потери полного столбцового ранга}
    \label{tab:results-check}
    \begin{tabular}{@{}lcc@{}}
        \toprule
        Вариант & $\delta_*$ & $\operatorname{rank} A'<n$ \\
        \midrule
        Регрессия, $\boldsymbol{A}_1$   & $0.075$      & $1<2$ \\
        Томография, $\boldsymbol{A}_1$ & $0.037037037$ & $1<2$ \\
        Томография, $\boldsymbol{A}_2$ & $0.1$         & $2<3$ \\
        \bottomrule
    \end{tabular}
\end{table}

Во всех случаях $\operatorname{rank}A'<n$, поэтому построенные точечные
матрицы имеют неполный столбцовый ранг. Для квадратной матрицы это
означает её вырожденность. Условие принадлежности
$|A'-A_c|\leq\delta_*D$ выполняется с погрешностью не более
$1.67\cdot10^{-16}$.

\subsection{Линейная регрессия: первый вариант}

Для интервальной матрицы $\boldsymbol{A}_1$ критический радиус и
соответствующий диапазон особенности равны
\[
    \delta_*=\frac{3}{40}=0.075,
    \qquad
    \delta\in[0.075,+\infty).
\]
Получены вектор и матрица
\[
    x_*=\begin{pmatrix}1\\-1.025\end{pmatrix},
    \qquad
    A'=\begin{pmatrix}
        1.025 & 1 \\
        1.025 & 1 \\
        1.025 & 1
    \end{pmatrix}.
\]
Все максимальные миноры $A'$ равны нулю.

\subsection{Малоракурсная томография: первая матрица}

Для интервальной матрицы $\boldsymbol{A}_1$ с радиусами всех элементов,
равными $\delta$, получено
\[
    \delta_*=\frac{1}{27}\approx0.037037037,
    \qquad
    \delta\in\left[\frac{1}{27},+\infty\right).
\]
Соответствующие вектор и матрица имеют вид
\[
    x_*=\begin{pmatrix}
        -0.4938271605\\
         0.5061728395
    \end{pmatrix},
    \qquad
    A'=\begin{pmatrix}
        0.987037037 & 0.962962963 \\
        1.037654321 & 1.012345679 \\
        1.062962963 & 1.037037037
    \end{pmatrix}.
\]
Все максимальные миноры $A'$ равны нулю с учётом погрешности вычислений.

\subsection{Малоракурсная томография: вторая матрица}

Для квадратной интервальной матрицы $\boldsymbol{A}_2$ получено
\[
    \delta_*=\frac{1}{10}=0.1,
    \qquad
    \delta\in[0.1,+\infty).
\]
Вектор и матрица на границе особенности равны
\[
    x_*=\begin{pmatrix}0\\-0.5\\0.5\end{pmatrix},
    \qquad
    A'=\begin{pmatrix}
        1.1 & 1.0 & 1.0 \\
        1.4 & 0.9 & 0.9 \\
        0.8 & 1.3 & 1.3
    \end{pmatrix}.
\]
Численно $\det A'=-2.02\cdot10^{-17}$, то есть с учётом погрешности
вычислений определитель равен нулю.

\section{Обсуждение}

\theorypart{Регрессия и томография}
Для одной и той же центральной матрицы $A_c$ получены разные критические
радиусы: $0.075$ в задаче регрессии и $1/27\approx0.037$ в задаче
томографии. В регрессии изменяется только первый столбец, а в томографии
--- все элементы, поэтому потеря ранга достигается при меньшем $\delta$.

\theorypart{Универсальность метода}
В основе всех расчётов использован один критерий
$|A_cx|\leq\delta D|x|$, не зависящий от квадратности матрицы.
Для регрессионного варианта из него получено аналитическое упрощение,
а для томографических вариантов применена общая оптимизационная схема.
Она без изменений позволила обработать матрицы размеров $3\times2$ и
$3\times3$, поэтому выбранный метод подходит как для прямоугольного,
так и для квадратного случая.

\section{Выводы}

\begin{enumerate}
    \item Реализован поиск минимального радиуса, при котором интервальная
    матрица содержит точечную матрицу неполного столбцового ранга.

    \item Для регрессионного варианта $\boldsymbol{A}_1$ получено
    $\delta_*=0.075$, для томографического варианта
    $\boldsymbol{A}_1$ --- $\delta_*=1/27$, для
    $\boldsymbol{A}_2$ --- $\delta_*=0.1$.

    \item Для каждого варианта построены точечная матрица $A'$ и
    ненулевой вектор $x_*$, удовлетворяющие условиям
    $A'\in\boldsymbol{A}(\delta_*)$ и $A'x_*=0$.

    \item Результаты показали, что критический радиус зависит не только
    от центральной матрицы, но и от структуры её допустимых возмущений.
\end{enumerate}

\section{Список литературы}
\begin{enumerate}
    \item Баженов~А.~Н., Карпова~А.~А.
    \emph{Интервальный анализ в примерах}: учебное пособие. ---
    Санкт-Петербург, 2024. --- 226~с.
\end{enumerate}

\section{Приложение}

Исходный код программы и электронная версия отчёта размещены в
GitHub-репозитории:
\url{https://github.com/kmozze/INTAN-lab1}.

\end{document}
